\documentclass[aspectratio=169]{beamer}

\usetheme{metropolis}
\usepackage[utf8]{inputenc}
\usepackage[T1]{fontenc}
\usepackage[brazil]{babel}
\usepackage{listings}
\usepackage{tabularx}
\usepackage{array}
\usepackage{booktabs}
\usepackage{tikz}
\usetikzlibrary{arrows.meta, positioning, shapes.geometric, calc, fit, backgrounds}

\setbeamertemplate{frame numbering}[fraction]
\metroset{block=fill}

\definecolor{termbg}{RGB}{40,44,52}
\definecolor{termfg}{RGB}{220,220,220}
\definecolor{codekw}{RGB}{235,129,27}
\definecolor{codecom}{RGB}{140,150,160}
\definecolor{codestr}{RGB}{152,195,121}

\lstdefinestyle{terminal}{
    backgroundcolor=\color{termbg},
    basicstyle=\ttfamily\footnotesize\color{termfg},
    breaklines=true,
    frame=single,
    rulecolor=\color{termbg},
    framerule=2pt,
    xleftmargin=4pt,
    xrightmargin=4pt,
    columns=fullflexible,
    keepspaces=true,
    showstringspaces=false
}

\lstdefinestyle{codigo}{
    style=terminal,
    language=C,
    keywordstyle=\color{codekw}\bfseries,
    commentstyle=\color{codecom}\itshape,
    stringstyle=\color{codestr},
    morekeywords={pid_t, size_t},
    moredelim=[is][\color{codekw}\bfseries]{@}{@}
}

% Elementos visuais das figuras
\tikzset{
    proc/.style={draw=mDarkTeal, fill=mDarkTeal!8, rounded corners=3pt, thick,
        minimum width=2.1cm, minimum height=0.8cm, align=center, font=\small},
    procf/.style={proc, draw=mDarkTeal!70, fill=white, minimum width=1.3cm, font=\footnotesize},
    destaque/.style={proc, draw=mLightBrown, fill=mLightBrown!12},
    arquivo/.style={draw=black!50, fill=black!5, thick, minimum width=2.2cm,
        minimum height=1.1cm, align=center, font=\small},
    seta/.style={-{Stealth[length=2.2mm]}, thick, draw=mDarkTeal},
    setad/.style={-{Stealth[length=2.2mm]}, very thick, draw=mLightBrown},
    rotulo/.style={font=\scriptsize, fill=white, inner sep=1.5pt},
    celula/.style={draw=mDarkTeal, minimum size=0.75cm, font=\ttfamily\small, fill=white}
}

% Revelar partes de uma figura TikZ aos poucos (overlays do beamer).
% visible on=<2->  : o elemento ocupa espaço sempre, mas só aparece a partir do slide 2;
% alt=<2>{a}{b}    : usa o estilo a no slide 2 e o estilo b nos demais.
\tikzset{
    invisivel/.style={opacity=0, text opacity=0},
    visible on/.style={alt={#1{}{invisivel}}},
    alt/.code args={<#1>#2#3}{\alt<#1>{\pgfkeysalso{#2}}{\pgfkeysalso{#3}}},
    novo/.style={draw=mLightBrown, fill=mLightBrown!25, very thick}
}

\title{Processos e comunicação entre processos}
\subtitle{fork, memória independente e pipe na prática}
\author{Weslei Ferreira Santos}
\institute{MATA58 · Sistemas Operacionais\\Instituto de Computação / UFBA}
\date{5 de outubro de 2026}

\begin{document}

\maketitle

\begin{frame}{Agenda}
\begin{center}
\small
\begin{tabular}{@{}r@{\hspace{1.2em}}l@{}}
\toprule
\textbf{min} & \textbf{Etapa} \\
\midrule
0 a 8   & Programa e processo \\
8 a 20  & Uma chamada de \texttt{fork()} \\
20 a 33 & \texttt{fork()} dentro do laço: \texttt{fork2.c} \\
33 a 45 & Memória independente e \texttt{waitpid} \\
45 a 53 & Checagem em duplas \\
53 a 65 & Pipe: canal, pontas e bloqueio \\
65 a 75 & Aplicação com vários processos: \texttt{pipe2.c} \\
75 a 87 & \textbf{Prática}: filho envia, pai recebe \\
87 a 90 & Bilhete de saída \\
\bottomrule
\end{tabular}
\end{center}
\end{frame}

\begin{frame}{A pergunta desta aula}
\begin{center}
{\Large \textbf{Como um processo recebe o resultado\\[0.2em] calculado por outro?}}
\end{center}

\vspace{1em}
\textbf{Ao final, você deve conseguir:}
\begin{itemize}
    \item identificar os retornos de \texttt{fork()} no pai e no filho;
    \item explicar por que cada processo tem a sua própria memória;
    \item desenhar a árvore de processos de um laço com \texttt{fork()};
    \item completar uma comunicação filho $\rightarrow$ pai com \texttt{pipe()}.
\end{itemize}
\end{frame}

% ============================================================
\section{Programa e processo}
% ============================================================

\begin{frame}{Um programa ou dois processos?}
\emph{Se eu executar o mesmo programa em dois terminais, tenho um programa ou dois processos?}

\vspace{0.8em}
\begin{center}
\begin{tikzpicture}
    \node[arquivo] (prog) {\texttt{fork2}\\\scriptsize arquivo no disco};
    \node[proc, right=3cm of prog, yshift=0.75cm] (p1) {Processo\\\scriptsize PID 4210};
    \node[proc, right=3cm of prog, yshift=-0.75cm] (p2) {Processo\\\scriptsize PID 4231};
    \draw[seta] (prog.east) -- node[rotulo, above, sloped] {executa} (p1.west);
    \draw[seta] (prog.east) -- node[rotulo, below, sloped] {executa} (p2.west);
    \node[right=0.4cm of p1, font=\scriptsize, align=left] {pilha, heap,\\registradores};
    \node[right=0.4cm of p2, font=\scriptsize, align=left] {pilha, heap,\\registradores};
\end{tikzpicture}
\end{center}

\pause
\begin{alertblock}{Definição}
O arquivo contém as instruções. Cada execução é um \textbf{processo}: uma instância com PID, registradores, pilha, heap e recursos próprios.
\end{alertblock}
\end{frame}

\begin{frame}{Pergunta para guardar}
\begin{center}
{\large Se uma dessas execuções altera uma variável,\\[0.3em]
a outra \emph{necessariamente} observa a alteração?}
\end{center}

\vspace{1.5em}
\centering Anote sua hipótese. Vamos testá-la com \texttt{fork2.c}.
\end{frame}

\begin{frame}{O que o SO guarda sobre cada processo}
Para cada processo, o núcleo mantém um registro (o \emph{bloco de controle do processo}):

\vspace{0.4em}
\begin{center}
\small
\begin{tabular}{@{}lp{9.6cm}@{}}
\toprule
\textbf{PID} & Número que identifica o processo enquanto ele existe. \\
\textbf{PPID} & PID do processo que o criou (o pai). \\
\textbf{Estado} & Executando, pronto, bloqueado esperando algo, terminado. \\
\textbf{Contexto} & Registradores, incluindo onde o programa parou; salvos quando ele sai da CPU. \\
\textbf{Memória} & O espaço de endereçamento: código, dados, heap e pilha. \\
\textbf{Descritores} & Arquivos, terminais e pipes que o processo tem abertos. \\
\bottomrule
\end{tabular}
\end{center}

\vspace{0.3em}
\small No terminal, \texttt{ps -f} mostra PID e PPID; \texttt{pstree -p} mostra quem criou quem.
\end{frame}

\begin{frame}{A memória de um processo}
\begin{columns}[c]
\column{0.42\textwidth}
\begin{tikzpicture}[node distance=0pt,
    regiao/.style={draw=mDarkTeal, fill=mDarkTeal!6, minimum width=3.4cm, minimum height=0.62cm, font=\small}]
    \node[regiao] (pilha) {pilha};
    \node[regiao, below=of pilha, fill=white, minimum height=0.9cm] (livre) {\scriptsize $\downarrow$ \quad livre \quad $\uparrow$};
    \node[regiao, below=of livre] (heap) {heap};
    \node[regiao, below=of heap] (dados) {dados globais};
    \node[regiao, below=of dados] (texto) {código};
\end{tikzpicture}

\column{0.56\textwidth}
\small
\begin{itemize}
    \item \textbf{Pilha}: variáveis locais e chamadas de função.
    \item \textbf{Heap}: \texttt{malloc}, \texttt{calloc}.
    \item \textbf{Dados}: variáveis globais e \texttt{static}.
    \item \textbf{Código}: as instruções do programa.
\end{itemize}
\end{columns}

\vspace{0.5em}
\begin{alertblock}{Endereços virtuais}
Os endereços que o programa usa são \textbf{virtuais}. O SO e o hardware traduzem cada um para uma posição física, e a tradução é \textbf{diferente para cada processo}. Visão simplificada: omite bibliotecas compartilhadas e o sorteio de endereços (ASLR).
\end{alertblock}
\end{frame}

% ============================================================
\section{Uma chamada de fork()}
% ============================================================

\begin{frame}[fragile]{Uma chamada, dois caminhos}
\begin{columns}[T]
\column{0.55\textwidth}
\begin{lstlisting}[style=codigo]
pid_t pid = fork();
if (pid == -1) {
    /* falha: nenhum filho criado */
} else if (pid == 0) {
    /* este caminho executa no filho */
} else {
    /* executa no pai; pid = filho */
}
/* pai e filho podem chegar aqui */
\end{lstlisting}

\column{0.42\textwidth}
\begin{center}
\begin{tabular}{@{}ll@{}}
\toprule
\textbf{Quem} & \textbf{fork() retorna} \\
\midrule
Filho & \texttt{0} \\
Pai   & PID do filho \\
Falha & \texttt{-1} \\
\bottomrule
\end{tabular}
\end{center}
\end{columns}

\vspace{0.6em}
\begin{alertblock}{O filho não recomeça o \texttt{main}}
Pai e filho continuam \textbf{depois} da chamada, cada um com seu retorno.
\end{alertblock}
\end{frame}

\begin{frame}{O que o filho recebe no \texttt{fork()}}
\begin{columns}[T]
\column{0.5\textwidth}
\textbf{Igual ao pai (cópia)}
\begin{itemize}
    \item Toda a memória: código, globais, heap, pilha.
    \item Valores das variáveis \emph{naquele instante}.
    \item Registradores: continua no mesmo ponto.
    \item Descritores abertos (arquivos, pipes).
\end{itemize}

\column{0.5\textwidth}
\textbf{Diferente do pai}
\begin{itemize}
    \item PID novo.
    \item PPID = PID do pai.
    \item Retorno de \texttt{fork()}: \texttt{0}.
\end{itemize}
\end{columns}

\vspace{0.6em}
\begin{exampleblock}{Cópia na escrita (\emph{copy-on-write})}
Pai e filho dividem as páginas até um deles escrever; só então a página é copiada. O efeito é o de uma cópia completa: \textbf{o que um escreve, o outro não vê}.
\end{exampleblock}
\end{frame}

\begin{frame}{\texttt{getpid()} e a ordem de execução}
\begin{itemize}
    \item \texttt{getpid()} retorna o PID do processo que \emph{está executando}.
    \item O valor que \texttt{fork()} devolve ao pai identifica o \textbf{filho}, não o próprio pai.
\end{itemize}

\vspace{1em}
\textbf{Previsão:} o pai necessariamente executa o próximo \texttt{printf} antes do filho?

\pause
\begin{exampleblock}{Resposta}
Não. A ordem depende do escalonamento e da sincronização que o programa estabelecer.
\end{exampleblock}
\end{frame}

\begin{frame}[fragile]{Vendo os dois retornos: \texttt{fork\_basico.c}}
\begin{lstlisting}[style=codigo, basicstyle=\ttfamily\scriptsize\color{termfg}]
pid_t pid = fork();
if (pid == 0) {
    printf("Filho: ... meu PID %d; meu pai %d\n", getpid(), getppid());
    fflush(stdout);  _exit(EXIT_SUCCESS);
}
printf("Pai: fork retornou %d; meu PID %d\n", pid, getpid());
\end{lstlisting}

\begin{lstlisting}[style=terminal, basicstyle=\ttfamily\scriptsize\color{termfg}]
$ make run-fork-basico
./build/fork-basico
Antes do fork: PID 4210
Filho: fork retornou 0; meu PID 4211; meu pai 4210
Pai: fork retornou 4211; meu PID 4210
Pai: filho 4211 terminou com codigo 0
\end{lstlisting}

\small O pai recebe 4211, que é o PID do filho; o pai do filho é 4210. Os PIDs mudam a cada execução.
\end{frame}

\begin{frame}{Um detalhe de C: \texttt{printf}, buffers e \texttt{\_exit}}
\begin{itemize}
    \item \texttt{printf} guarda o texto num \textbf{buffer na memória do processo} e só entrega ao SO depois (quebra de linha no terminal, buffer cheio ou fim do programa).
    \item O buffer faz parte da memória, então \textbf{é copiado pelo \texttt{fork()}}: texto pendente pode sair duas vezes.
    \item \texttt{\_exit} encerra \textbf{sem esvaziar} o buffer. \texttt{exit} e \texttt{return} esvaziam, inclusive o que foi herdado do pai.
\end{itemize}

\vspace{0.5em}
\begin{alertblock}{Regra prática}
\texttt{fflush(stdout)} antes do \texttt{fork()}. No filho que termina no meio do código do pai, \texttt{fflush(stdout)} e depois \texttt{\_exit}.
\end{alertblock}
\end{frame}

% ============================================================
\section{fork() dentro do laço}
% ============================================================

\begin{frame}[fragile]{\texttt{fork2.c}: o trecho central}
\begin{lstlisting}[style=codigo]
int *v = calloc((size_t)tamanho, sizeof *v);   /* tamanho = 3 */
pid_t original = getpid();
for (int i = 0; i < tamanho; ++i) {
    pid_t pid = fork();
    if (pid == -1) { /* tratamento de falha */ }
    if (pid == 0) v[i] = i + 5;
}
\end{lstlisting}

\vspace{0.5em}
\begin{alertblock}{Dois minutos, no papel}
Com \texttt{tamanho = 3}, quantos processos existem ao todo? Desenhe quem cria quem.
\end{alertblock}
\end{frame}

\begin{frame}{Lendo \texttt{fork2.c} por partes}
\small
\begin{center}
\begin{tabular}{@{}>{\raggedright\arraybackslash}p{4.9cm}>{\raggedright\arraybackslash}p{8.1cm}@{}}
\toprule
\textbf{Trecho} & \textbf{O que faz} \\
\midrule
\texttt{calloc(tamanho, ...)} & Reserva o vetor \textbf{já zerado}: \texttt{[0, 0, 0]}. \\
\texttt{original = getpid()} & Guarda o PID de quem começou. Os filhos herdam o número, mas só um processo tem \texttt{getpid() == original}. \\
laço com \texttt{fork()} & Em cada volta, \textbf{todo} processo vivo chama \texttt{fork()}. O filho escreve \texttt{v[i] = i + 5}. \\
laço com \texttt{waitpid} & Cada processo espera \textbf{os seus} filhos diretos, até não restar nenhum (\texttt{ECHILD}). \\
\texttt{if (getpid()==original)} & Só o processo original soma o vetor e imprime. \\
\bottomrule
\end{tabular}
\end{center}

\vspace{0.3em}
\centering Pergunta-chave: \textbf{o filho sai do laço depois de nascer?} Não: ele continua as voltas que faltam.
\end{frame}

\begin{frame}{A árvore de \texttt{fork2.c}, volta a volta}
\begin{columns}[T]
\column{0.62\textwidth}
\begin{tikzpicture}[x=1.15cm, y=1.05cm, procf/.append style={minimum width=0.8cm},
    vet/.style={font=\tiny\ttfamily, text=mDarkTeal, inner sep=1pt}]
    % Etapa 1: antes do laço, só o original
    \node[proc, minimum width=1.6cm] (p0) at (0,0) {original};
    % Etapa 2 (i = 0): o original cria A
    \node[procf, visible on=<2->, alt=<2>{novo}{}] (a) at (2,0) {A};
    \draw[seta, visible on=<2->] (p0) -- node[rotulo, above, visible on=<2->] {i=0} (a);
    % Etapa 3 (i = 1): original cria B; A cria C
    \node[procf, visible on=<3->, alt=<3>{novo}{}] (c) at (3.5,0) {C};
    \node[procf, visible on=<3->, alt=<3>{novo}{}] (b) at (2,-2) {B};
    \draw[seta, visible on=<3->] (a) -- node[rotulo, above, visible on=<3->] {i=1} (c);
    \draw[seta, visible on=<3->] (p0.south) |- node[rotulo, pos=0.8, above, visible on=<3->] {i=1} (b.west);
    % Etapa 4 (i = 2): original cria D; A cria E; B cria F; C cria G
    \node[procf, visible on=<4->, alt=<4>{novo}{}] (g) at (5,0) {G};
    \node[procf, visible on=<4->, alt=<4>{novo}{}] (e) at (3.5,-1) {E};
    \node[procf, visible on=<4->, alt=<4>{novo}{}] (f) at (3.5,-2) {F};
    \node[procf, visible on=<4->, alt=<4>{novo}{}] (d) at (2,-3) {D};
    \draw[seta, visible on=<4->] (c) -- node[rotulo, above, visible on=<4->] {i=2} (g);
    \draw[seta, visible on=<4->] (a.south) |- node[rotulo, pos=0.75, above, visible on=<4->] {i=2} (e.west);
    \draw[seta, visible on=<4->] (b) -- node[rotulo, above, visible on=<4->] {i=2} (f);
    \draw[seta, visible on=<4->] (p0.south) ++(-0.3,0) |- node[rotulo, pos=0.8, above, visible on=<4->] {i=2} (d.west);
    % Etapa 5: o vetor v de cada processo ao final
    \node[vet, above=1pt of p0, visible on=<5>] {[0,0,0]};
    \node[vet, above=1pt of a, visible on=<5>] {[5,0,0]};
    \node[vet, above=1pt of c, visible on=<5>] {[5,6,0]};
    \node[vet, above=1pt of b, visible on=<5>] {[0,6,0]};
    \node[vet, right=2pt of g, visible on=<5>] {[5,6,7]};
    \node[vet, right=2pt of e, visible on=<5>] {[5,0,7]};
    \node[vet, right=2pt of f, visible on=<5>] {[0,6,7]};
    \node[vet, right=2pt of d, visible on=<5>] {[0,0,7]};
\end{tikzpicture}

\vspace{0.2em}
{\scriptsize \textcolor{mLightBrown}{\textbf{laranja}}: processos que nasceram nesta volta}

\column{0.36\textwidth}
\textbf{Processos ao fim de cada volta}
\begin{itemize}
    \item<1-> início: \textbf{1}
    \item<2-> \texttt{i=0}: \textbf{2} \hfill {\scriptsize +A}
    \item<3-> \texttt{i=1}: \textbf{4} \hfill {\scriptsize +B, C}
    \item<4-> \texttt{i=2}: \textbf{8} \hfill {\scriptsize +D, E, F, G}
\end{itemize}

\vspace{0.4em}
\small
\only<1>{Antes do laço existe só o processo \textbf{original}.}%
\only<2>{O original chama \texttt{fork()} e nasce \textbf{A}. A escreve \texttt{v[0] = 5} na \emph{sua} cópia.}%
\only<3>{Original \textbf{e A} chamam \texttt{fork()}: A também continua o laço! Nascem \textbf{B} e \textbf{C}.}%
\only<4>{Os quatro chamam \texttt{fork()}, em qualquer ordem. Nascem \textbf{D, E, F e G}.}%
\only<5>{\textbf{8 processos}, 7 novos. O código tem \textbf{uma} chamada de \texttt{fork()}; a árvore executou $1+2+4=7$.\\[0.3em]
Ao lado de cada um, o seu vetor \texttt{v}: o 18 só existe em G.}
\end{columns}
\end{frame}

\begin{frame}{Dois cuidados com essa contagem}
\begin{itemize}
    \item \textbf{8} conta os processos criados \emph{ao longo da execução}.\\
    Não significa que os oito estejam vivos ao mesmo tempo.
    \item O código tem \textbf{uma} chamada de \texttt{fork()}; a árvore executa \textbf{sete}.
\end{itemize}

\vspace{1em}
\textbf{E se cada filho saísse do laço logo após nascer?}

\pause
\begin{exampleblock}{Resposta}
Só o original completaria as três iterações: \textbf{4 processos} ao todo (o original e três filhos).
\end{exampleblock}
\end{frame}

% ============================================================
\section{Memória independente}
% ============================================================

\begin{frame}[fragile]{Previsão antes de executar}
Cada filho escreve em \texttt{v[i]}. O pai soma o vetor no final.

\begin{center}
{\large Qual soma o pai imprime? \quad \textbf{0}? \quad \textbf{18}? \quad outro valor?}
\end{center}

\pause
\begin{lstlisting}[style=terminal]
$ make run-fork
./build/fork2
Soma total no pai: 0
\end{lstlisting}
\end{frame}

\begin{frame}{Cada processo escreve na própria cópia}
\begin{center}
\begin{tikzpicture}
    \node[font=\small\bfseries, anchor=east] at (-0.2,1.6) {Pai antes};
    \foreach \x in {0,1,2} \node[celula] at (\x*0.75+0.4,1.6) {0};
    \draw[seta] (2.4,1.6) -- node[rotulo, above] {fork} (4.2,2.4);
    \draw[seta] (2.4,1.6) -- (4.2,0.8);

    \node[font=\small\bfseries, anchor=west] at (4.3,2.4) {Filho};
    \node[celula, fill=mLightBrown!25] at (6.0,2.4) {5};
    \node[celula] at (6.75,2.4) {0};
    \node[celula] at (7.5,2.4) {0};
    \node[font=\scriptsize, anchor=west] at (8.0,2.4) {\texttt{v[0] = 0 + 5}};

    \node[font=\small\bfseries, anchor=west] at (4.3,0.8) {Pai};
    \foreach \x in {0,1,2} \node[celula] at (\x*0.75+6.0,0.8) {0};
    \node[font=\scriptsize, anchor=west] at (8.0,0.8) {não entra em \texttt{pid == 0}};
\end{tikzpicture}
\end{center}

\begin{itemize}
    \item \texttt{calloc} começa com três zeros.
    \item O filho altera a \textbf{sua} memória; o vetor do pai continua \texttt{[0, 0, 0]}.
    \item \texttt{getpid() == original}: só o processo original imprime a soma.
\end{itemize}
\end{frame}

\begin{frame}{O vetor de cada um dos oito processos}
\begin{columns}[c]
\column{0.5\textwidth}
\small
\begin{tabular}{@{}lllc@{}}
\toprule
\textbf{Proc.} & \textbf{Criado por} & \textbf{Vetor} & \textbf{Soma} \\
\midrule
original & início & \texttt{[0,0,0]} & 0 \\
A & original, i=0 & \texttt{[5,0,0]} & 5 \\
B & original, i=1 & \texttt{[0,6,0]} & 6 \\
D & original, i=2 & \texttt{[0,0,7]} & 7 \\
C & A, i=1 & \texttt{[5,6,0]} & 11 \\
E & A, i=2 & \texttt{[5,0,7]} & 12 \\
F & B, i=2 & \texttt{[0,6,7]} & 13 \\
G & C, i=2 & \texttt{[5,6,7]} & \textbf{18} \\
\bottomrule
\end{tabular}

\column{0.48\textwidth}
\begin{itemize}
    \item O filho herda o vetor \textbf{como estava no pai} no momento do \texttt{fork()}.
    \item Depois escreve só a sua posição.
    \item O 18 existe, mas só em \textbf{G}, que não é o original e não imprime.
\end{itemize}
\end{columns}
\end{frame}

\begin{frame}[fragile]{Mesmo endereço, valores diferentes}
\begin{lstlisting}[style=codigo]
int x = 10;
pid_t pid = fork();
if (pid == 0) { x = 99; printf("Filho: x = %d em %p\n", x, (void *)&x); ... }
printf("Pai:   x = %d em %p\n", x, (void *)&x);   /* depois de esperar o filho */
\end{lstlisting}

\begin{lstlisting}[style=terminal]
$ make run-memoria
./build/memoria-independente
Filho: x = 99 em 0x7ffe72feea88
Pai:   x = 10 em 0x7ffe72feea88
\end{lstlisting}

\begin{alertblock}{Mesmo endereço não é mesma memória}
O endereço é \textbf{virtual}: cada processo o interpreta no seu próprio espaço, e o SO o leva a posições físicas diferentes.
\end{alertblock}
\end{frame}

\begin{frame}{\texttt{waitpid}: esperar não é comunicar}
\begin{columns}[T]
\column{0.5\textwidth}
\textbf{O que \texttt{waitpid} faz}
\begin{itemize}
    \item Espera um filho terminar.
    \item Coleta o estado de término.
    \item Cada processo coleta seus \textbf{filhos diretos}.
\end{itemize}

\column{0.5\textwidth}
\textbf{O que ele não faz}
\begin{itemize}
    \item Não copia o vetor do filho.
    \item Não cria memória compartilhada.
    \item Não transporta resultados.
\end{itemize}
\end{columns}

\vspace{1em}
\begin{alertblock}{Por que a soma não virou 18?}
O pai já espera os filhos, mas esperar coordena o \emph{término}. Nem \texttt{sleep} nem \texttt{waitpid} tornam o heap compartilhado.
\end{alertblock}
\end{frame}

\begin{frame}[fragile]{\texttt{waitpid} na prática}
\begin{lstlisting}[style=codigo]
pid_t waitpid(pid_t pid, int *status, int opcoes);
\end{lstlisting}

\begin{columns}[T]
\column{0.48\textwidth}
\small
\begin{itemize}
    \item \texttt{pid > 0}: espera esse filho; \texttt{-1}: qualquer filho.
    \item \texttt{opcoes = 0}: \textbf{bloqueia} até um filho terminar.
    \item Devolve o PID coletado ou \texttt{-1}.
    \item \texttt{EINTR}: interrompido por sinal, chame de novo.
    \item \texttt{ECHILD}: não há mais filhos.
\end{itemize}

\column{0.5\textwidth}
\small
\begin{tabular}{@{}ll@{}}
\toprule
\textbf{Macro} & \textbf{Responde} \\
\midrule
\texttt{WIFEXITED(st)} & Terminou normalmente? \\
\texttt{WEXITSTATUS(st)} & Com que código? \\
\texttt{WIFSIGNALED(st)} & Morto por sinal? \\
\texttt{WTERMSIG(st)} & Qual sinal? \\
\bottomrule
\end{tabular}
\end{columns}
\end{frame}

\begin{frame}{Zumbis e órfãos}
\begin{center}
\begin{tikzpicture}[node distance=1.3cm]
    \node[proc] (exec) {executando};
    \node[destaque, right=of exec] (zumbi) {zumbi};
    \node[proc, right=of zumbi, fill=white] (fim) {removido};
    \draw[seta] (exec) -- node[rotulo, above] {termina} (zumbi);
    \draw[seta] (zumbi) -- node[rotulo, above] {pai coleta} (fim);
\end{tikzpicture}
\end{center}

\begin{columns}[T]
\column{0.5\textwidth}
\textbf{Zumbi}\\
\small Terminou, mas o pai ainda não coletou. Fica na tabela de processos (estado \texttt{Z} no \texttt{ps}) guardando o código de saída.

\column{0.5\textwidth}
\textbf{Órfão}\\
\small O pai terminou antes. Outro processo adota e coleta: o PID 1 num servidor, ou o \texttt{systemd -{}-user} numa sessão gráfica.
\end{columns}

\vspace{0.6em}
\begin{alertblock}{Por que coletar?}
Um pai que nunca chama \texttt{waitpid} acumula zumbis.
\end{alertblock}
\end{frame}

\begin{frame}[fragile]{E se o filho devolvesse o 30 no código de saída?}
\begin{lstlisting}[style=codigo]
_exit(30);                     /* no filho */
WEXITSTATUS(status) == 30      /* no pai: funciona! */
\end{lstlisting}

\pause
\textbf{Mas não serve como canal de dados:}
\begin{itemize}
    \item Só \textbf{8 bits}: valores de 0 a 255. \texttt{\_exit(300)} chega como \texttt{44}.
    \item Não cabe número negativo, \texttt{double}, texto ou vetor.
    \item Só chega quando o filho \textbf{termina}: um único valor, no fim.
\end{itemize}

\vspace{0.4em}
\begin{exampleblock}{Para que serve}
Dizer \emph{como} o processo terminou: \texttt{0} é sucesso, outro valor indica o tipo de falha.
\end{exampleblock}
\end{frame}

\begin{frame}{Checagem em duplas \hfill \normalsize 5 minutos}
\begin{enumerate}
    \item Quantos processos são criados ao todo, incluindo o original? Desenhe a evolução por iteração.
    \item Quais retornos de \texttt{fork()} identificam filho, pai e falha? De onde cada um continua?
    \item O filho altera \texttt{v[i]}. Qual soma o pai imprime? Explique.
    \item O pai espera os filhos. Isso faz o vetor dele receber as escritas dos filhos?
\end{enumerate}

\vspace{0.6em}
\centering \emph{Vamos ouvir uma dupla sobre a árvore e outra sobre a espera.}
\end{frame}

% ============================================================
\section{Pipe: comunicação explícita}
% ============================================================

\begin{frame}{Se a memória não leva o resultado\ldots}
\begin{center}
{\large \ldots o filho precisa \textbf{enviar os dados explicitamente}.}
\end{center}

\vspace{0.4em}
\textbf{IPC} (\emph{Inter-Process Communication}): mecanismos do SO para processos trocarem dados.

\vspace{0.6em}
\begin{center}
\begin{tikzpicture}
    \node[proc] (filho) {Filho\\\scriptsize calcula 30};
    \node[draw=mLightBrown, fill=mLightBrown!12, thick, minimum width=3cm, minimum height=0.8cm,
          right=2.8cm of filho, font=\small] (pipe) {pipe};
    \node[proc, right=2.8cm of pipe] (pai) {Pai\\\scriptsize recebe 30};
    \draw[setad] (filho) -- node[rotulo, above=2pt] {\texttt{write(fd[1])}} (pipe);
    \draw[setad] (pipe) -- node[rotulo, above=2pt] {\texttt{read(fd[0])}} (pai);
    \node[below=0.15cm of filho, font=\scriptsize] {fecha \texttt{fd[0]}};
    \node[below=0.15cm of pai, font=\scriptsize] {fecha \texttt{fd[1]}};
    \begin{scope}[on background layer]
        \node[draw=black!40, dashed, rounded corners, fit=(pipe), inner sep=0.35cm,
              label={[font=\scriptsize, text=black!60]below:mantido pelo núcleo do SO}] {};
    \end{scope}
\end{tikzpicture}
\end{center}
\end{frame}

\begin{frame}{Pontas, descritores e herança}
\begin{itemize}
    \item \texttt{pipe(fd)} cria o canal: \texttt{fd[0]} é a ponta de \textbf{leitura}, \texttt{fd[1]} a de \textbf{escrita}.
    \item \textbf{Descritor}: número que identifica um recurso aberto pelo processo.
    \item O pipe é criado \textbf{antes} do \texttt{fork()}: pai e filho herdam descritores que referenciam o mesmo canal.
    \item Cada processo fecha a ponta que não usa.
\end{itemize}

\vspace{0.6em}
\begin{alertblock}{Ordem que importa}
\texttt{pipe()} $\rightarrow$ \texttt{fork()} $\rightarrow$ cada um fecha uma ponta $\rightarrow$ \texttt{write} / \texttt{read} $\rightarrow$ \texttt{waitpid}
\end{alertblock}
\end{frame}

\begin{frame}{A tabela de descritores depois do \texttt{fork()}}
\begin{center}
\begin{tikzpicture}[
    tab/.style={draw=mDarkTeal, minimum width=1.9cm, minimum height=0.5cm, font=\ttfamily\scriptsize, fill=white},
    tabx/.style={tab, fill=black!8, text=black!45}]
    \node[font=\small\bfseries] at (0,2.25) {Filho};
    \node[tab] (f0) at (0,1.7) {0 entrada};
    \node[tab, below=0pt of f0] (f1) {1 saída};
    \node[tab, below=0pt of f1] (f2) {2 erros};
    \node[tabx, below=0pt of f2] (f3) {3 fechado};
    \node[tab, below=0pt of f3, fill=mLightBrown!15] (f4) {4 = fd[1]};
    \node[font=\small\bfseries] at (9,2.25) {Pai};
    \node[tab] (p0) at (9,1.7) {0 entrada};
    \node[tab, below=0pt of p0] (p1) {1 saída};
    \node[tab, below=0pt of p1] (p2) {2 erros};
    \node[tab, below=0pt of p2, fill=mLightBrown!15] (p3) {3 = fd[0]};
    \node[tabx, below=0pt of p3] (p4) {4 fechado};
    \node[draw=mLightBrown, fill=mLightBrown!12, thick, minimum width=3cm, minimum height=0.7cm, font=\small]
        (pipe) at (4.5,-0.2) {buffer do pipe};
    \node[font=\scriptsize, color=black!60, below=0.05cm of pipe] {núcleo do SO};
    \draw[setad] (f4.east) -- node[rotulo, above, sloped] {escreve} (pipe.west);
    \draw[setad] (pipe.east) -- node[rotulo, above, sloped] {lê} (p3.west);
\end{tikzpicture}
\end{center}

\small
\begin{itemize}
    \item \texttt{0}, \texttt{1} e \texttt{2} já vêm abertos; por isso o pipe costuma receber \texttt{3} e \texttt{4}.
    \item O \texttt{fork()} copia a tabela: as duas cópias apontam para o \textbf{mesmo} pipe do núcleo.
    \item Pipe criado \textbf{depois} do \texttt{fork()}: cada um teria o seu, sem canal em comum.
\end{itemize}
\end{frame}

\begin{frame}[fragile]{O exemplo pequeno: \texttt{pipe\_soma.c}}
\begin{lstlisting}[style=codigo]
if (pid == 0) {                                   /* filho */
    close(fd[0]);
    valor = 10 + 20;
    int resultado = transferir(fd[1], &valor, sizeof valor, 1);
    close(fd[1]);
    _exit(resultado == 0 ? EXIT_SUCCESS : EXIT_FAILURE);
}
close(fd[1]);                                     /* pai */
int resultado = transferir(fd[0], &valor, sizeof valor, 0);
close(fd[0]);
\end{lstlisting}

\small \texttt{transferir} chama \texttt{write} (modo 1) ou \texttt{read} (modo 0) até mover todos os bytes.
\end{frame}

\begin{frame}[fragile]{Por dentro do helper \texttt{transferir}}
\begin{lstlisting}[style=codigo]
while (tamanho > 0) {
    ssize_t n = modo ? write(fd, p, tamanho) : read(fd, p, tamanho);
    if (n < 0 && errno == EINTR) continue;    /* interrompido: tenta de novo */
    if (n < 0) return -1;                     /* erro de verdade */
    if (n == 0) return -1;                    /* canal fechado antes do fim */
    p += n;                                   /* avanca no dado */
    tamanho -= (size_t)n;                     /* quanto ainda falta */
}
\end{lstlisting}

\small
\begin{itemize}
    \item \texttt{read} e \texttt{write} podem mover \textbf{menos bytes} do que o pedido; por isso o laço.
    \item Na prática vocês \textbf{usam} o helper; o foco é quem fecha qual ponta, quem envia e quem recebe.
\end{itemize}
\end{frame}

\begin{frame}{O mesmo \texttt{\&valor}, duas variáveis}
\begin{itemize}
    \item \texttt{\&valor} no filho e \texttt{\&valor} no pai apontam para variáveis \textbf{locais de cada processo}.
    \item O SO \textbf{copia os bytes} pelo canal.
    \item O ponteiro não torna as duas variáveis compartilhadas: o endereço é interpretado dentro do espaço de cada processo.
\end{itemize}

\vspace{0.6em}
\begin{exampleblock}{Pipe é fluxo de bytes}
Nem toda chamada de \texttt{read} recebe tudo que pedimos. O \emph{helper} \texttt{transferir} repete a operação e trata interrupções (\texttt{EINTR}).
\end{exampleblock}
\end{frame}

\begin{frame}{E se o pai chegar ao \texttt{read} primeiro?}
\begin{center}
\begin{tabular}{@{}p{6.2cm}p{6.2cm}@{}}
\toprule
\textbf{Situação do pipe} & \textbf{O que \texttt{read} faz} \\
\midrule
Vazio, com ponta de escrita aberta & Bloqueia e espera dados \\
Com dados disponíveis & Retorna os bytes (talvez menos que o pedido) \\
Vazio, todas as escritas fechadas & Retorna \texttt{0}: \textbf{EOF} \\
\bottomrule
\end{tabular}
\end{center}

\vspace{0.8em}
\begin{alertblock}{Por isso fechamos pontas}
Enquanto houver \emph{qualquer} ponta de escrita aberta, inclusive uma cópia herdada esquecida, um pipe vazio não sinaliza EOF.
\end{alertblock}
\end{frame}

\begin{frame}{E do lado de quem escreve?}
\begin{center}
\small
\begin{tabular}{@{}p{5.6cm}p{7.2cm}@{}}
\toprule
\textbf{Situação} & \textbf{O que \texttt{write} faz} \\
\midrule
Há espaço no buffer do pipe & Copia os bytes e retorna. \\
Buffer cheio (64\,KiB no Linux) & \textbf{Bloqueia} até o leitor consumir. \\
Nenhuma ponta de leitura aberta & Sinal \texttt{SIGPIPE}: encerra o processo. Se ignorado, \texttt{-1} com \texttt{EPIPE}. \\
\bottomrule
\end{tabular}
\end{center}

\vspace{0.5em}
\begin{itemize}
    \item Escritas de até \texttt{PIPE\_BUF} bytes (4096 no Linux) são \textbf{atômicas}: não se misturam com as de outro processo.
    \item Um \texttt{int} cabe com folga.
    \item \texttt{pipe2.c} ignora \texttt{SIGPIPE} para poder tratar o erro em vez de morrer.
\end{itemize}
\end{frame}

\begin{frame}[fragile]{Executando o exemplo}
\begin{lstlisting}[style=terminal]
$ make run-pipe-soma
./build/pipe-soma
Resultado recebido: 30
\end{lstlisting}

\vspace{0.6em}
\begin{itemize}
    \item O \texttt{30} foi calculado no filho.
    \item Chegou ao pai pelo canal, não por memória comum.
\end{itemize}
\end{frame}

% ============================================================
\section{Vários processos: pipe2.c}
% ============================================================

\begin{frame}{\texttt{pipe2.c}: quem cria quem}
\begin{center}
\begin{tikzpicture}[node distance=0.6cm and 1.2cm]
    \node[destaque] (pai) {Pai original};
    \node[proc, below left=1cm and 1.2cm of pai] (calc) {Calculador\\\scriptsize\texttt{son1}};
    \node[proc, below right=1cm and 1.2cm of pai] (coord) {Coordenador\\\scriptsize\texttt{son5}};
    \node[procf, minimum width=2cm, below=1cm of coord] (f2) {fornece \texttt{2a}};
    \node[procf, minimum width=2cm, left=0.4cm of f2] (f1) {fornece \texttt{-b}};
    \node[procf, minimum width=2cm, right=0.4cm of f2] (f3) {fornece \texttt{delta}};
    \draw[seta] (pai) -- (calc);
    \draw[seta] (pai) -- (coord);
    \draw[seta] (coord) -- (f1);
    \draw[seta] (coord) -- (f2);
    \draw[seta] (coord) -- (f3);
\end{tikzpicture}
\end{center}

\vspace{0.3em}
\centering \textbf{Seis processos} e \textbf{cinco pipes}. A entrada é lida no pai antes dos \texttt{fork()}.
\end{frame}

\begin{frame}{\texttt{pipe2.c}: por onde os dados passam}
\begin{center}
\begin{tikzpicture}[node distance=0.35cm]
    \node[procf, minimum width=2.1cm] (f1) {\texttt{-b}};
    \node[procf, minimum width=2.1cm, below=of f1] (f2) {\texttt{2a}};
    \node[procf, minimum width=2.1cm, below=of f2] (f3) {\texttt{delta}};
    \node[proc, right=3.2cm of f2, minimum height=1.6cm] (calc) {Calculador};
    \node[destaque, right=3.2cm of calc] (pai) {Pai};
    \draw[setad] (f1.east) -- node[rotulo, above, sloped] {\texttt{canais[0]}} (calc.west);
    \draw[setad] (f2.east) -- node[rotulo, above] {\texttt{canais[1]}} (calc.west);
    \draw[setad] (f3.east) -- node[rotulo, below, sloped] {\texttt{canais[2]}} (calc.west);
    \draw[setad] ([yshift=0.25cm]calc.east) -- node[rotulo, above] {\texttt{canais[3]}: x1} ([yshift=0.25cm]pai.west);
    \draw[setad] ([yshift=-0.25cm]calc.east) -- node[rotulo, below] {\texttt{canais[4]}: x2} ([yshift=-0.25cm]pai.west);
\end{tikzpicture}
\end{center}

\vspace{0.3em}
\begin{alertblock}{Árvore de criação $\neq$ fluxo de dados}
Os fornecedores são filhos do coordenador, mas enviam ao calculador. Quem escreve e quem lê cada canal?
\end{alertblock}
\end{frame}

\begin{frame}{\texttt{pipe2.c}: cada processo fica só com as pontas que usa}
\begin{center}
\small
\begin{tabular}{@{}lll@{}}
\toprule
\textbf{Processo} & \textbf{Lê de} & \textbf{Escreve em} \\
\midrule
Pai & \texttt{canais[3]}, \texttt{canais[4]} & nenhum \\
Calculador & \texttt{canais[0]}, \texttt{[1]}, \texttt{[2]} & \texttt{canais[3]}, \texttt{[4]} \\
Coordenador & nenhum & \texttt{canais[0]}, \texttt{[1]}, \texttt{[2]}, só até criar os fornecedores \\
Fornecedor \texttt{i} & nenhum & \texttt{canais[i]} \\
\bottomrule
\end{tabular}
\end{center}

\small \texttt{fechar\_exceto} fecha todas as pontas, menos as indicadas. O coordenador guarda as de escrita só para os fornecedores herdarem, e logo depois fecha.

\begin{alertblock}{Por que tanto cuidado?}
Uma ponta de escrita esquecida aberta impede o EOF: o calculador poderia esperar para sempre.
\end{alertblock}
\end{frame}

\begin{frame}[fragile]{Executando \texttt{pipe2.c}}
\begin{lstlisting}[style=terminal]
$ make run-pipe
./build/pipe2
>>> Digite tres numeros separados por espaco e tecle Enter: a b delta
>>> Exemplo: 1 -3 1   (delta ja calculado; o programa nao pede c)
a b delta: 1 -3 1
Raizes: x1 = 1; x2 = 2
\end{lstlisting}

\begin{columns}[T]
\column{0.5\textwidth}
\textbf{Entrada:} \texttt{a}, \texttt{b} e \texttt{delta}.\\
O programa \textbf{não} pede \texttt{c}.

\column{0.5\textwidth}
$-b = 3, \quad 2a = 2, \quad \sqrt{\Delta} = 1$\\[0.3em]
$x_1 = \frac{3-1}{2} = 1, \quad x_2 = \frac{3+1}{2} = 2$
\end{columns}

\vspace{0.6em}
\small O exemplo ilustra IPC; não foi feito para acelerar um cálculo tão pequeno.
\end{frame}

% ============================================================
\section{Prática}
% ============================================================

\begin{frame}[fragile]{Seis lacunas: \texttt{pipe\_soma\_atividade.c}}
\begin{lstlisting}[style=codigo]
if (pid == 0) {
    close(@/* 1: ponta que o filho nao usa */@);
    valor = @/* 2: calculo */@;
    int resultado = @/* 3: enviar valor pelo helper transferir */@;
    close(fd[1]);
    _exit(resultado == 0 ? EXIT_SUCCESS : EXIT_FAILURE);
}
close(@/* 4: ponta que o pai nao usa */@);
int resultado = @/* 5: receber valor pelo helper transferir */@;
close(@/* 6: ponta usada pelo pai */@);
\end{lstlisting}

\small \texttt{transferir(descritor, endereço, bytes, modo)}: modo \texttt{1} envia, modo \texttt{0} recebe.
\end{frame}

\begin{frame}{Antes de preencher: perguntas que guiam cada lacuna}
\begin{columns}[T]
\column{0.5\textwidth}
\textbf{No filho}
\begin{enumerate}
    \item O filho só \textbf{escreve}. Qual ponta ele nunca usa?
    \item Qual conta ele precisa fazer?
    \item Envie \texttt{valor}: qual descritor, qual endereço, quantos bytes, qual modo?
\end{enumerate}

\column{0.5\textwidth}
\textbf{No pai}
\begin{enumerate}
    \setcounter{enumi}{3}
    \item O pai só \textbf{lê}. Qual ponta ele nunca usa?
    \item Receba em \texttt{valor}: qual descritor, endereço, bytes e modo?
    \item Terminou de ler: qual ponta falta fechar?
\end{enumerate}
\end{columns}

\vspace{0.6em}
\begin{alertblock}{Confira}
Endereço é \texttt{\&valor} (não \texttt{valor}); tamanho é \texttt{sizeof valor}; \texttt{fd[0]} lê, \texttt{fd[1]} escreve.
\end{alertblock}
\end{frame}

\begin{frame}[fragile]{Compilar e conferir}
\begin{lstlisting}[style=terminal]
$ gcc -std=c11 -Wall -Wextra -Wpedantic -Werror \
      examples/dia05/pipe_soma_atividade.c -o /tmp/pipe-soma-aluno
$ /tmp/pipe-soma-aluno
Resultado recebido: 30
\end{lstlisting}

\begin{columns}[T]
\column{0.5\textwidth}
\textbf{Em duplas, 7 minutos}
\begin{itemize}
    \item Antes de preencher, o esqueleto não compila.
    \item Cada pessoa explica uma ponta do canal.
\end{itemize}

\column{0.5\textwidth}
\textbf{Responda também}
\begin{enumerate}
    \setcounter{enumi}{4}
    \item Por que criar o pipe antes do \texttt{fork()}?
    \item \texttt{read} antes do \texttt{write}: erro, espera ou EOF?
    \item Por que cada um fecha uma ponta diferente?
\end{enumerate}
\end{columns}
\end{frame}

% ============================================================
\section{Fechamento}
% ============================================================

\begin{frame}{Erros comuns}
\begin{center}
\small
\begin{tabular}{@{}p{4.6cm}p{8.4cm}@{}}
\toprule
\textbf{Se você pensou\ldots} & \textbf{Reveja\ldots} \\
\midrule
``O pai soma 18.'' & Em qual processo cada atribuição acontece? Memórias separadas. \\
``São quatro processos.'' & O filho não sai do laço: ele também chama \texttt{fork()}. \\
``\texttt{waitpid} compartilha memória.'' & Esperar o término não transporta bytes. \\
``\texttt{read} sempre recebe tudo.'' & O retorno diz quantos bytes chegaram. \\
``Pipe vazio sempre dá EOF.'' & Ainda há alguma ponta de escrita aberta? \\
``Mesmo endereço, mesma memória.'' & O endereço é virtual, interpretado em cada processo. \\
``O filho recomeça o \texttt{main}.'' & Ele continua logo depois do \texttt{fork()}. \\
\bottomrule
\end{tabular}
\end{center}
\end{frame}

\begin{frame}{Síntese}
\begin{center}
\begin{tabular}{@{}lp{9.5cm}@{}}
\toprule
\texttt{fork()}    & Cria um processo com memória \textbf{independente}. \\
\texttt{waitpid()} & Coordena o \textbf{término} e coleta o filho. \\
\texttt{pipe()}    & \textbf{Transporta dados} entre processos. \\
\bottomrule
\end{tabular}
\end{center}

\vspace{0.8em}
\begin{exampleblock}{Próxima aula: 7/10}
Threads \textbf{compartilham} memória. O problema muda: como controlar o acesso a ela?
\end{exampleblock}
\end{frame}

\begin{frame}{Bilhete de saída}
\begin{center}
{\large \textbf{Individual, em duas ou três frases:}}

\vspace{1em}
{\Large Por que \texttt{waitpid} não resolve a comunicação\\[0.2em] do resultado e \texttt{pipe} resolve?}
\end{center}
\end{frame}

\begin{frame}{Glossário}
\small
\begin{columns}[T]
\column{0.5\textwidth}
\textbf{Processo}: programa em execução, com memória e recursos próprios.\\[0.35em]
\textbf{PID / PPID}: número do processo / do seu pai.\\[0.35em]
\textbf{\texttt{fork()}}: cria um filho, cópia do pai.\\[0.35em]
\textbf{\texttt{waitpid()}}: espera um filho terminar e o coleta.\\[0.35em]
\textbf{Zumbi}: terminou, mas ainda não foi coletado.\\[0.35em]
\textbf{Órfão}: o pai terminou antes; outro processo adota.

\column{0.5\textwidth}
\textbf{IPC}: comunicação entre processos.\\[0.35em]
\textbf{Pipe}: canal de bytes de mão única, mantido pelo SO.\\[0.35em]
\textbf{Descritor}: número que o processo usa para um recurso aberto.\\[0.35em]
\textbf{\texttt{fd[0]} / \texttt{fd[1]}}: ponta de leitura / de escrita.\\[0.35em]
\textbf{Bloquear}: o processo dorme até a condição mudar.\\[0.35em]
\textbf{EOF}: \texttt{read} devolve 0; não há mais quem escreva.
\end{columns}
\end{frame}

\begin{frame}[fragile]{Para estudar depois}
\textbf{Rode os exemplos} (na raiz do repositório):
\begin{lstlisting}[style=terminal]
make run-fork-basico    make run-fork    make run-memoria
make run-pipe-soma      make run-pipe    # digite 1 -3 1
\end{lstlisting}

\textbf{Leia}
\begin{itemize}
    \item A \textbf{apostila} da aula: os mesmos temas, explicados com mais calma, e exercícios com respostas.
    \item \texttt{man 2 fork}, \texttt{man 2 waitpid}, \texttt{man 2 pipe}, \texttt{man 7 pipe}.
    \item Tanenbaum e Bos, \emph{Sistemas Operacionais Modernos}, cap. 2.
    \item Silberschatz, Galvin e Gagne, \emph{Fundamentos de Sistemas Operacionais}, cap. 3.
\end{itemize}
\end{frame}

\begin{frame}{Obrigado!}
\centering
{\Large Dúvidas?}
\vspace{2em}

Weslei Ferreira Santos\\
MATA58 · Sistemas Operacionais\\
Instituto de Computação / UFBA

\vspace{1.5em}
{\scriptsize Material-base: Prof. Maycon Leone M. Peixoto}
\end{frame}

\end{document}
